\documentclass[11pt,a4paper]{article}

\usepackage[utf8]{inputenc}
\usepackage[T1]{fontenc}
\usepackage{lmodern}
\usepackage[margin=2.5cm]{geometry}
\usepackage{amsmath,amssymb,amsthm}
\usepackage{bm}
\usepackage{booktabs}
\usepackage{hyperref}
\usepackage{siunitx}

\newcommand{\jj}{\mathrm{j}}
\newcommand{\ee}{\mathrm{e}}
\newcommand{\diff}{\mathrm{d}}
\newcommand{\matZ}{\bm{Z}}
\newcommand{\matY}{\bm{Y}}
\newcommand{\matR}{\bm{R}}
\newcommand{\matL}{\bm{L}}
\newcommand{\matC}{\bm{C}}
\newcommand{\matG}{\bm{G}}
\newcommand{\matA}{\bm{A}}
\newcommand{\matB}{\bm{B}}
\newcommand{\matD}{\bm{D}}
\newcommand{\matM}{\bm{M}}
\newcommand{\matF}{\bm{F}}
\newcommand{\matP}{\bm{P}}
\newcommand{\matT}{\bm{T}}
\newcommand{\matI}{\bm{I}}
\newcommand{\vV}{\bm{V}}
\newcommand{\vI}{\bm{I}}
\newcommand{\vQ}{\bm{Q}}
\newcommand{\vphi}{\bm{\phi}}
\newcommand{\vu}{\bm{u}}
\newcommand{\vp}{\bm{p}}
\newcommand{\vq}{\bm{q}}
\newcommand{\zeroM}{\bm{0}}
\newcommand{\Lint}{L_{\mathrm{int}}}
\newcommand{\Lself}[1]{L^{\mathrm{geom}}_{\mathrm{self},#1}}
\newcommand{\Lmut}[1]{L^{\mathrm{geom}}_{\mathrm{mut},#1}}
\renewcommand{\Re}{\operatorname{Re}}

\title{Distributed-Parameter Modelling of a Three-Phase\\
Overhead Transmission Line and Phase-Current\\
Unbalance as a Function of Length and Apparent Power}
\author{Litgrid}
\date{}

\begin{document}
\maketitle

\begin{abstract}
This note documents the analytical model used by the Litgrid application to
compute voltage and current profiles along a three-phase overhead transmission
line. Per-unit-length matrices \(\matR,\matL,\matC,\matG\) are obtained from
tower geometry by decomposing each conductor's self impedance into an internal
(skin) contribution and an external (geometric) contribution, building a
fully-coupled cable-to-cable matrix, reducing it first across the bundled
conductors of a phase, and finally Kron-eliminating the overhead ground wires.
The line is then solved as a distributed-parameter system using the matrix
exponential of its state-space generator. A two-point boundary problem is
posed with a symmetrical positive-sequence voltage at one terminal and a
specified three-phase complex power \(S=P+\jj Q\) at the same (or opposite)
terminal. The natural asymmetry of the per-unit-length matrices for
non-transposed geometries drives the phase-current unbalance, which is
quantified and presented as a function of length and apparent power. A
complete worked example for a \SI{110}{\kilo\volt} single-circuit line closes
the document.
\end{abstract}

\section{Notation and Conventions}

Bold symbols denote vectors and matrices in the conductor- or phase-domain.
Phasors are RMS, with the cosine reference. The phase rotation operator is
\(a = \ee^{\jj 2\pi/3}\) and the positive-sequence unit vector is
\begin{equation}
\vu \;=\; \begin{bmatrix} 1 \\ a^{2} \\ a \end{bmatrix}.
\label{eq:u}
\end{equation}
Angular frequency \(\omega=2\pi f\). Distances along the line are in
kilometres, so \(\matR\) is in \si{\ohm/\kilo\metre}, \(\matL\) in
\si{\milli\henry/\kilo\metre}, \(\matC\) in \si{\nano\farad/\kilo\metre}, and
\(\matG\) in \si{\micro\siemens/\kilo\metre}.

\section{Per-Unit-Length Parameters of the Coupled Conductor System}

The tower carries \(N\) conductors in total, comprising the bundled
sub-conductors of \(n_p=3\) phase bundles plus \(n_g\ge 0\) overhead ground
wires. Each conductor \(i\) is described by its DC resistance
\(r_{\mathrm{dc},i}\) at the operating temperature, its outer radius \(r_i\),
and its Cartesian position \((x_i,y_i)\) with \(y_i>0\) above the perfectly
conducting earth plane.

\subsection{Self and Mutual Inductances}

Following the classical decomposition of the magnetic flux linking a thin
round non-magnetic conductor at the power frequency, the self inductance per
unit length of conductor \(i\) is the sum of two contributions:

\paragraph{Internal inductance.} The flux that resides \emph{inside} the
conductor (uniform-current low-frequency approximation):
\begin{equation}
\Lint \;=\; \frac{\mu_{0}}{8\pi}\quad[\si{\henry/\metre}].
\label{eq:Lint}
\end{equation}
This is the same for every solid round conductor and does not depend on its
radius.

\paragraph{Geometric (external) self inductance.} The flux that links the
conductor through the region between its surface (radius \(r_i\)) and the
earth-return reference distance \(D_e\) below the soil:
\begin{equation}
\Lself{i}\;=\;\frac{\mu_{0}}{2\pi}\,\ln\!\frac{D_{e}}{r_{i}}
\quad[\si{\henry/\metre}].
\label{eq:Lselfgeom}
\end{equation}

\paragraph{Geometric mutual inductance.} Between two distinct conductors
\(i\) and \(j\) separated by the (real) distance
\(d_{ij}=\sqrt{(x_{i}-x_{j})^{2}+(y_{i}-y_{j})^{2}}\):
\begin{equation}
\Lmut{ij}\;=\;\frac{\mu_{0}}{2\pi}\,\ln\!\frac{D_{e}}{d_{ij}}
\quad[\si{\henry/\metre}].
\label{eq:Lmut}
\end{equation}

The earth-return depth and the earth-return resistance follow Carson's
low-frequency homogeneous-soil approximation:
\begin{equation}
D_{e} \;=\; 658.5\,\sqrt{\rho/f}\;\;[\si{\metre}],\qquad
R_{e} \;=\; \frac{\omega\mu_{0}}{8\pi}\;\;[\si{\ohm/\metre}],
\label{eq:CarsonDeRe}
\end{equation}
where \(\rho\) is the soil resistivity in \si{\ohm.\metre}.

\subsection{Series Impedance Matrix of the Coupled System}

Assembling \eqref{eq:Lint}--\eqref{eq:Lmut} into the per-unit-length series
impedance gives, in \si{\ohm/\kilo\metre},
\begin{align}
Z_{ii} &= \bigl(r_{\mathrm{dc},i} + R_{e}\bigr)
   \;+\; \jj\,\omega\,\bigl(\Lint + \Lself{i}\bigr)\cdot 10^{3}, \label{eq:Zii}\\[3pt]
Z_{ij} &= R_{e}
   \;+\; \jj\,\omega\,\Lmut{ij}\cdot 10^{3},\qquad i\ne j. \label{eq:Zij}
\end{align}
The \(N\times N\) symmetric matrix \(\matZ_{\mathrm{full}}\) collects every
cable-to-cable coupling (phase sub-conductors and ground wires alike).

\subsection{Potential Coefficient Matrix of the Coupled System}

With the earth treated as a perfect equipotential, image-charge theory yields
the \(N\times N\) potential-coefficient matrix \(\matP_{\mathrm{full}}\):
\begin{align}
P_{ii} &= \frac{1}{2\pi\varepsilon_{0}}\,
          \ln\!\frac{2y_{i}}{r_{i}}, \label{eq:Pii}\\[3pt]
P_{ij} &= \frac{1}{2\pi\varepsilon_{0}}\,
          \ln\!\frac{D_{ij}}{d_{ij}},\qquad i\ne j, \label{eq:Pij}
\end{align}
where \(D_{ij}=\sqrt{(x_{i}-x_{j})^{2}+(y_{i}+y_{j})^{2}}\) is the distance
from conductor \(i\) to the \emph{image} of conductor \(j\) below the soil.
\(\matP_{\mathrm{full}}\) relates charge to voltage:
\(\vV_{\mathrm{full}}=\matP_{\mathrm{full}}\,\vQ_{\mathrm{full}}\).

\section{Reduction to the Phase Domain}

The \(N\)-conductor coupled matrices contain information that is not directly
useful for line operation: sub-conductors of a bundle are paralleled, and the
overhead ground wires are short-circuited to earth. Two successive linear
reductions eliminate these degrees of freedom.

\subsection{Bundle Reduction}

Let the \(N\) physical conductors be partitioned into \(n_p\) phase bundles
plus \(n_g\) ground wires (each ground wire is its own ``bundle of size one'').
Introduce the connection matrix \(\matT\in\{0,1\}^{N\times(n_{p}+n_{g})}\)
with \(T_{ik}=1\) if conductor \(i\) belongs to bundle \(k\), zero otherwise.
The two physical constraints of a bundle are
\begin{equation}
\vV_{\mathrm{full}}=\matT\,\vV_{b},\qquad
\vI_{b}=\matT^{\top}\vI_{\mathrm{full}},\qquad
\vQ_{b}=\matT^{\top}\vQ_{\mathrm{full}},
\label{eq:bundle-T}
\end{equation}
expressing equality of potential within a bundle and Kirchhoff summation of
the per-conductor currents and charges.

Substituting in \(\vV=\matZ_{\mathrm{full}}\vI\) and
\(\vV=\matP_{\mathrm{full}}\vQ\) and using
\(\matY_{\mathrm{full}}=\matZ_{\mathrm{full}}^{-1}\) yields the bundle-reduced
matrices
\begin{align}
\matZ_{b} &= \bigl(\matT^{\top}\matZ_{\mathrm{full}}^{-1}\,\matT\bigr)^{-1},
\label{eq:bundle-Z}\\[3pt]
\matP_{b} &= \bigl(\matT^{\top}\matP_{\mathrm{full}}^{-1}\,\matT\bigr)^{-1}.
\label{eq:bundle-P}
\end{align}
Both have size \((n_{p}+n_{g})\times(n_{p}+n_{g})\). When every bundle is a
single conductor, \(\matT=\matI_{N}\) and the bundle step is the identity.

\subsection{Kron Elimination of the Ground Wires}

Partition the bundle-reduced matrices into a phase block (subscript \(p\),
size \(n_{p}=3\)) and a ground-wire block (subscript \(g\), size \(n_{g}\)):
\begin{equation}
\matZ_{b}=\begin{bmatrix}\matZ_{pp}&\matZ_{pg}\\[2pt]\matZ_{gp}&\matZ_{gg}\end{bmatrix},\qquad
\matP_{b}=\begin{bmatrix}\matP_{pp}&\matP_{pg}\\[2pt]\matP_{gp}&\matP_{gg}\end{bmatrix}.
\end{equation}
Imposing \(\vV_{g}=\bm{0}\) on the ground wires (perfectly bonded to earth)
removes the ground-block degrees of freedom and gives the \(3\times 3\)
phase-domain matrices
\begin{align}
\matZ_{\mathrm{ph}} &= \matZ_{pp} - \matZ_{pg}\,\matZ_{gg}^{-1}\,\matZ_{gp},
\label{eq:Kron-Z}\\[2pt]
\matP_{\mathrm{ph}} &= \matP_{pp} - \matP_{pg}\,\matP_{gg}^{-1}\,\matP_{gp}.
\label{eq:Kron-P}
\end{align}
The phase-domain per-unit-length matrices follow as
\begin{equation}
\matR = \Re\{\matZ_{\mathrm{ph}}\},\quad
\matL = \tfrac{1}{\omega}\,\mathrm{Im}\{\matZ_{\mathrm{ph}}\}\cdot 10^{3},\quad
\matC = \matP_{\mathrm{ph}}^{-1}\cdot 10^{12},\quad
\matG = g_{0}\,\matI_{3},
\label{eq:RLCG}
\end{equation}
with \(g_{0}\) a small leakage conductance modelling corona/insulator losses
(default \SI{2e-3}{\micro\siemens/\kilo\metre}). For non-transposed geometries
\(\matZ_{\mathrm{ph}}\) and \(\matP_{\mathrm{ph}}\) are symmetric but
\emph{not} circulant; their unequal entries are the structural origin of
every unbalance phenomenon that follows.

\section{Distributed-Parameter Line Equations}

Let \(\vV(s)\) and \(\vI(s)\) denote the phasor voltage and current vectors
at distance \(s\in[0,\ell]\) from the local terminal. Form the per-unit-length
series impedance and shunt admittance
\begin{equation}
\matZ \;=\; \matR + \jj\omega\,\matL,\qquad
\matY \;=\; \matG + \jj\omega\,\matC.
\end{equation}
The Telegrapher's equations in the frequency domain are
\begin{equation}
\frac{\diff}{\diff s}\!
\begin{bmatrix}\vV\\[2pt]\vI\end{bmatrix}
\;=\;
\underbrace{\begin{bmatrix}\zeroM & -\matZ \\[2pt] -\matY & \zeroM\end{bmatrix}}_{\matF}
\begin{bmatrix}\vV\\[2pt]\vI\end{bmatrix}.
\label{eq:state-space}
\end{equation}
The state vector at any distance follows from the matrix exponential
\begin{equation}
\vphi(s)\;\equiv\;\begin{bmatrix}\vV(s)\\[2pt]\vI(s)\end{bmatrix}
\;=\;\ee^{\matF s}\,\vphi(0).
\label{eq:propagator}
\end{equation}
The end-to-end ABCD transfer is the \(6\times 6\) matrix
\(\matM=\ee^{\matF\ell}\), partitioned as
\begin{equation}
\matM=\begin{bmatrix}\matA&\matB\\\matA'&\matD\end{bmatrix},\qquad
\begin{bmatrix}\vV_{r}\\[2pt]\vI_{r}\end{bmatrix}
=\matM\begin{bmatrix}\vV_{s}\\[2pt]\vI_{s}\end{bmatrix},
\label{eq:ABCD}
\end{equation}
with subscripts \(s\) and \(r\) denoting the sending end (\(s=0\)) and the
receiving end (\(s=\ell\)).

\section{Boundary Conditions: Symmetric Voltage and Specified \(S\)}

Two physical constraints close the boundary problem:

\begin{enumerate}
\item \textbf{Balanced symmetric positive-sequence voltage at the local
terminal.} The local bus is the angular reference and is held at the nominal
line-to-neutral magnitude \(V_{s}=V_{LL}/\sqrt{3}\):
\begin{equation}
\vV_{s}\;=\;V_{s}\,\vu.
\label{eq:bc-V}
\end{equation}

\item \textbf{Specified three-phase complex power at the local terminal.}
\begin{equation}
S\;=\;P+\jj Q\;=\;\vV_{s}^{\!\top}\,\vI_{s}^{*}.
\label{eq:bc-S}
\end{equation}
(An entirely symmetric formulation applies if the constraint is given at the
remote terminal: swap \(s\leftrightarrow r\), invert \(\matM\), and proceed
identically.)
\end{enumerate}

The remote bus is assumed sufficiently stiff to remain balanced in magnitude
and angle, so
\begin{equation}
\vV_{r}\;=\;V_{r}\,\ee^{\jj\delta}\,\vu,
\label{eq:Vr-ansatz}
\end{equation}
where \(V_{r}\) and \(\delta\) are two unknown real scalars. Define the
single complex unknown
\begin{equation}
\boxed{\;W\;\equiv\;V_{r}\,\ee^{-\jj\delta}\;\;\Longleftrightarrow\;\;V_{r}=|W|,\;\;\delta=-\arg(W),\;}
\label{eq:W-def}
\end{equation}
so that \(V_{r}\,\ee^{\jj\delta}=W^{*}\).

\section{Closed-Form Solution for \(W\)}

Solving the upper block of \eqref{eq:ABCD} for \(\vI_{s}\) and substituting
\eqref{eq:bc-V}--\eqref{eq:Vr-ansatz} gives
\begin{equation}
\vI_{s}\;=\;\matB^{-1}\bigl(\vV_{r}-\matA\,\vV_{s}\bigr)
\;=\;\matB^{-1}\bigl(W^{*}\vu - V_{s}\,\matA\,\vu\bigr).
\label{eq:Is-from-bc}
\end{equation}
Plugging \eqref{eq:Is-from-bc} into \eqref{eq:bc-S} yields a single scalar
complex equation linear in \(W\). Expanding and collecting terms,
\begin{equation}
S\;=\;V_{s}\,W\,\underbrace{\vu\cdot\overline{\matB^{-1}\vu}}_{\alpha}
\;-\;V_{s}^{2}\,\underbrace{\vu\cdot\overline{\matB^{-1}\matA\,\vu}}_{\beta},
\label{eq:S-linear-in-W}
\end{equation}
where the dot products use ordinary (non-Hermitian) component sums and the
overline denotes complex conjugation. Both \(\alpha,\beta\in\mathbb{C}\) depend
only on the line (through \(\matA\) and \(\matB\)), so
\begin{equation}
\boxed{\;
W\;=\;\frac{S\;+\;V_{s}^{2}\,\beta}{V_{s}\,\alpha},
\qquad V_{r}=|W|,\quad \delta=-\arg(W).
\;}
\label{eq:W-closed-form}
\end{equation}
Once \(W\) is known, \eqref{eq:Is-from-bc} delivers \(\vI_{s}\); the receiving-end voltage and current follow from \eqref{eq:Vr-ansatz} and \eqref{eq:ABCD}.

\section{Analytical Form of the Phase-Current Magnitudes}
\label{sec:analytical-Ik}

Substituting \(\vV_{s}=V_{s}\vu\) and \(\vV_{r}=W^{*}\vu\) in
\eqref{eq:Is-from-bc} and extracting the \(k\)-th component yields a
remarkably compact expression. Introduce the two complex 3-vectors
\begin{equation}
\vp \;:=\; \matB^{-1}\vu,\qquad
\vq \;:=\; \matB^{-1}\matA\,\vu,
\label{eq:pq-def}
\end{equation}
both of which depend on the line geometry only (through
\(\matM=\ee^{\matF\ell}\)) and are independent of \(W\) and \(S\). Then
\begin{equation}
I_{k} \;=\; W^{*}\,p_{k} \;-\; V_{s}\,q_{k},\qquad k\in\{A,B,C\}.
\label{eq:Ik-linear-W}
\end{equation}

\subsection{Closed Form for \texorpdfstring{\(|I_{k}|\)}{|Ik|}}

Taking modulus squared of \eqref{eq:Ik-linear-W} and using
\(W=V_{r}\ee^{-\jj\delta}\),
\begin{equation}
|I_{k}|^{2}
\;=\; V_{r}^{2}\,|p_{k}|^{2}
\;+\; V_{s}^{2}\,|q_{k}|^{2}
\;-\; 2\,V_{s} V_{r}\,|p_{k}|\,|q_{k}|\,\cos\!\bigl(\delta + \psi_{k}\bigr),
\label{eq:Ik-cosrule}
\end{equation}
where the \(\psi_{k}\) are line-only phase angles
\begin{equation}
\psi_{k} \;:=\; \arg(p_{k}) - \arg(q_{k}),\qquad k\in\{A,B,C\}.
\label{eq:psi-def}
\end{equation}
Equation \eqref{eq:Ik-cosrule} is the law of cosines: each \(|I_{k}|\) is the
length of the third side of a triangle with sides \(V_{r}|p_{k}|\) and
\(V_{s}|q_{k}|\) enclosing the angle \((\delta+\psi_{k})\).

\subsection{Largest Modular Deviation Between Phases}

The peak-to-peak spread of the three phase magnitudes is, in closed form,
\begin{equation}
\boxed{\;
\Delta_{\max}(W) \;=\;
\max_{k}\!\sqrt{\,V_{r}^{2}|p_{k}|^{2}+V_{s}^{2}|q_{k}|^{2}-2V_{s}V_{r}|p_{k}||q_{k}|\cos(\delta+\psi_{k})\,}
\;-\;\min_{k}\!\sqrt{\,(\text{same})\,}.\;}
\label{eq:Deltamax}
\end{equation}
Equivalently, every pairwise modulus difference reads
\begin{equation}
|I_{i}|^{2}-|I_{j}|^{2}
\;=\;V_{r}^{2}\bigl(|p_{i}|^{2}-|p_{j}|^{2}\bigr)
\;+\;V_{s}^{2}\bigl(|q_{i}|^{2}-|q_{j}|^{2}\bigr)
\;-\;2V_{s}V_{r}\bigl[|p_{i}||q_{i}|\cos(\delta+\psi_{i})-|p_{j}||q_{j}|\cos(\delta+\psi_{j})\bigr],
\label{eq:pairwise-diff}
\end{equation}
which exposes the three sources of asymmetry:
\begin{itemize}
\item the spread of the moduli \(|p_{k}|\), weighted by \(V_{r}^{2}\)
(receiving-side, shunt-driven asymmetry);
\item the spread of the moduli \(|q_{k}|\), weighted by \(V_{s}^{2}\)
(sending-side, series-driven asymmetry);
\item a cross term modulated by \(\cos(\delta+\psi_{k})\), which is the only
piece that depends on the transmission angle \(\delta\) and hence is the
only piece sensitive to the loading direction.
\end{itemize}
The relative deviation reported by the application
(eq.~\ref{eq:devs}) follows by normalising \eqref{eq:Ik-cosrule} by the mean
of the three magnitudes; the maximum percentage deviation is therefore an
explicit function of \(W\) alone, given the six line constants
\(\{|p_{k}|,|q_{k}|,\psi_{k}\}_{k\in\{A,B,C\}}\).

\section{Voltage and Current Profile Along the Line}

With the initial state \(\vphi(0)=[\vV_{s};\,\vI_{s}]\) determined, the entire
profile is recovered from \eqref{eq:propagator}:
\begin{equation}
\vphi(s_{k})\;=\;\ee^{\matF s_{k}}\,\vphi(0),
\qquad s_{k}\in\{0,\Delta s,\dots,\ell\}.
\end{equation}
Each phase voltage and phase current is plotted as a function of \(s\). The
receiving-end power is recomputed by direct dot product
\(S_{r}=\vV_{r}^{\!\top}\vI_{r}^{*}\), and the total active loss is
\(P_{\mathrm{loss}}=P-\Re\{S_{r}\}\).

\section{Quantifying Phase-Current Unbalance}

\subsection{Per-Phase Magnitude and Angular-Spacing Deviations}

At a chosen terminal, define magnitudes \(|I_{k}|\) and angles
\(\theta_{k}=\arg(I_{k})\) for \(k\in\{A,B,C\}\), the mean magnitude
\(\overline{|I|}=\tfrac{1}{3}(|I_{A}|+|I_{B}|+|I_{C}|)\), and the three signed
inter-phase angular spacings wrapped to \((-\pi,\pi]\):
\begin{equation}
s_{AB}=\mathrm{wrap}(\theta_{B}-\theta_{A}),\;
s_{BC}=\mathrm{wrap}(\theta_{C}-\theta_{B}),\;
s_{CA}=\mathrm{wrap}(\theta_{A}-\theta_{C}),
\end{equation}
with mean absolute spacing
\(\overline{|s|}=\tfrac{1}{3}(|s_{AB}|+|s_{BC}|+|s_{CA}|)\).
The per-phase deviations reported are
\begin{equation}
\Delta_{\mathrm{mag},k}[\%]=\frac{|I_{k}|-\overline{|I|}}{\overline{|I|}}\cdot 100,\qquad
\Delta_{\mathrm{ang},k}[\si{\degree}]=|s_{k}|-\overline{|s|}.
\label{eq:devs}
\end{equation}

\subsection{Spatial Variation Along the Line}

Because the per-unit-length matrices are non-circulant, each phase current
magnitude is not constant along \(s\):
\begin{equation}
\eta_{k}[\%]\;=\;\frac{\max_{s}|I_{k}(s)|-\min_{s}|I_{k}(s)|}{\min_{s}|I_{k}(s)|}\cdot 100,
\quad k\in\{A,B,C\}.
\label{eq:eta}
\end{equation}

\subsection{Relation to the Standard Current Unbalance Factor}

For comparison with international standards, the symmetrical-component decomposition
\begin{equation}
\begin{bmatrix}I_{0}\\I_{+}\\I_{-}\end{bmatrix}=\tfrac{1}{3}
\begin{bmatrix}1&1&1\\1&a&a^{2}\\1&a^{2}&a\end{bmatrix}
\begin{bmatrix}I_{A}\\I_{B}\\I_{C}\end{bmatrix}
\end{equation}
yields the current unbalance factor
\begin{equation}
\mathrm{IUF}\;=\;\frac{|I_{-}|}{|I_{+}|}\cdot 100\;[\%].
\label{eq:IUF}
\end{equation}

\section{Unbalance as a Function of Length and Apparent Power}

Equations \eqref{eq:W-closed-form}--\eqref{eq:Is-from-bc} make the dependence
of the solution on \(\ell\) and \(S\) explicit. The line-dependent quantities
\(\matA(\ell)\) and \(\matB(\ell)\) come exclusively from
\(\matM=\ee^{\matF\ell}\), so:
\begin{itemize}
\item In the electrically short limit \(\matA\to\matI_{3}\) and
\(\matB\to\ell\,\matZ\), so the sending current reduces to
\(\vI_{s}\approx(\ell\,\matZ)^{-1}(W^{*}\vu - V_{s}\vu)\): the dominant
asymmetry is that of \(\matZ\) itself, scaled by \(1/\ell\).
\item As \(\ell\) grows the shunt admittance contributes through
\(\matB^{-1}\matA\); the asymmetry inherited from the non-circulant geometry
propagates into both \(\alpha\) and \(\beta\).
\item For fixed line geometry, \(S=P+\jj Q\) sets the absolute magnitude of
\(\vI_{s}\) but only weakly modifies the relative unbalance
\eqref{eq:devs}--\eqref{eq:eta}, because the constant capacitive term
\(V_{s}^{2}\beta\) becomes a smaller fraction of the loading term as \(|S|\)
grows. Light loading therefore typically produces a larger percentage
unbalance than heavy loading.
\end{itemize}
The Litgrid application sweeps \((\ell,S)\) by re-evaluating
\eqref{eq:W-closed-form}, \eqref{eq:Is-from-bc} and
\eqref{eq:devs}--\eqref{eq:eta} for each pair.

%========================================================================
\section{Worked Example: \texorpdfstring{\SI{110}{\kilo\volt}}{110 kV} Single-Circuit Line}
%========================================================================

\subsection{Inputs}

A horizontal-vertical asymmetric configuration with one conductor per phase
and one overhead ground wire.
\begin{itemize}
\item Nominal line-to-line voltage \(V_{LL}=\SI{110}{\kilo\volt}\), so the
local line-to-neutral magnitude is \(V_{s}=110/\sqrt{3}\approx
\SI{63509.1}{\volt}\).
\item Three-phase complex power injected at the local terminal:
\(S=P+\jj Q=80+\jj 10\;\si{\mega\volt\ampere}\).
\item Frequency \(f=\SI{50}{\hertz}\), so \(\omega=100\pi
\approx\SI{314.159}{\radian/\second}\).
\item Soil resistivity \(\rho=\SI{1000}{\ohm.\metre}\).
\item Phase conductors (one per phase, no bundling): outer diameter
\(\SI{17.1}{\milli\metre}\) (so \(r_{p}=\SI{8.55}{\milli\metre}\)),
\(r_{\mathrm{dc}}=\SI{0.1897}{\ohm/\kilo\metre}\).
\item Ground wire: outer diameter \(\SI{11.11}{\milli\metre}\)
(\(r_{g}=\SI{5.555}{\milli\metre}\)),
\(r_{\mathrm{dc}}=\SI{2.53}{\ohm/\kilo\metre}\).
\item Geometry (order \(A,B,C,G\)):
\(x=(-4,\,0,\,+4,\,0)\,\si{\metre}\) and
\(y=(20,\,20,\,23,\,25)\,\si{\metre}\).
\end{itemize}
There is one conductor per phase, so the bundle reduction is the identity
(\(\matT=\matI_{4}\)). Only the ground-wire elimination is non-trivial.

\subsection{Geometric Distances and Image Distances}

Real distances \(d_{ij}\,[\si{\metre}]\) and image distances
\(D_{ij}\,[\si{\metre}]\):
\begin{center}
\begin{tabular}{lcccc}
\toprule
pair & $A$--$B$ & $B$--$C$ & $A$--$C$ & $A$--$G$ \\
$d_{ij}$ & $4.000$ & $5.000$ & $\sqrt{73}\!\approx\!8.544$ & $\sqrt{41}\!\approx\!6.403$ \\
$D_{ij}$ & $\sqrt{1616}\!\approx\!40.199$ & $\sqrt{1865}\!\approx\!43.186$ & $\sqrt{1913}\!\approx\!43.738$ & $\sqrt{2041}\!\approx\!45.177$ \\
\midrule
pair & $B$--$G$ & $C$--$G$ & & \\
$d_{ij}$ & $5.000$ & $\sqrt{20}\!\approx\!4.472$ & & \\
$D_{ij}$ & $45.000$ & $\sqrt{2320}\!\approx\!48.166$ & & \\
\bottomrule
\end{tabular}
\end{center}
The self-image distances are \(2y_{A}=2y_{B}=40\), \(2y_{C}=46\),
\(2y_{G}=50\;\si{\metre}\).

\subsection{Carson Constants}

\begin{align*}
D_{e} &= 658.5\sqrt{1000/50} \;\approx\; \SI{2944.65}{\metre},\\
R_{e} &= \frac{\omega\mu_{0}}{8\pi}\cdot 10^{3}
        \;\approx\; \SI{0.01571}{\ohm/\kilo\metre},\\
\frac{\omega\mu_{0}}{2\pi}\cdot 10^{3}
      &\approx\; \SI{0.06283}{\ohm/\kilo\metre\;per\;unit\;\ln}.
\end{align*}
Note that numerically the internal reactance per kilometre equals the
earth-return resistance: \(\omega\,\Lint\cdot 10^{3}=R_{e}=\SI{0.01571}{\ohm/\kilo\metre}\).

\subsection{Coupled \texorpdfstring{$4\times 4$}{4x4} Series Impedance Matrix}

Apply \eqref{eq:Zii}--\eqref{eq:Zij}. For instance,
\begin{align*}
Z_{AA} &= (0.1897 + 0.01571) + \jj\bigl[0.01571 + 0.06283\,\ln(2944.65/0.00855)\bigr]\\
       &= 0.2054 + \jj\,0.8167 \;[\si{\ohm/\kilo\metre}],\\[3pt]
Z_{AB} &= 0.01571 + \jj\,0.06283\,\ln(2944.65/4) \;=\; 0.01571 + \jj\,0.4149,\\
Z_{GG} &= (2.53 + 0.01571) + \jj\bigl[0.01571 + 0.06283\,\ln(2944.65/0.005555)\bigr]\\
       &= 2.5457 + \jj\,0.8438.
\end{align*}
Repeating for every pair, the fully coupled cable-to-cable matrix
\(\matZ_{\mathrm{full}}\), in \si{\ohm/\kilo\metre}, with rows/columns ordered
\((A,B,C,G)\), is
\begin{equation*}
\matZ_{\mathrm{full}}=
\begin{bmatrix}
0.2054{+}\jj 0.8167 & 0.0157{+}\jj 0.4149 & 0.0157{+}\jj 0.3671 & 0.0157{+}\jj 0.3852\\
0.0157{+}\jj 0.4149 & 0.2054{+}\jj 0.8167 & 0.0157{+}\jj 0.4008 & 0.0157{+}\jj 0.4008\\
0.0157{+}\jj 0.3671 & 0.0157{+}\jj 0.4008 & 0.2054{+}\jj 0.8167 & 0.0157{+}\jj 0.4078\\
0.0157{+}\jj 0.3852 & 0.0157{+}\jj 0.4008 & 0.0157{+}\jj 0.4078 & 2.5457{+}\jj 0.8438
\end{bmatrix}.
\end{equation*}

\subsection{Kron Elimination of the Ground Wire}

With \(\matT=\matI_{4}\), the bundle-reduced matrix coincides with
\(\matZ_{\mathrm{full}}\). Partitioning the last row/column as the ground
block and applying \eqref{eq:Kron-Z} produces the \(3\times 3\) phase-domain
matrix
\begin{equation*}
\matZ_{\mathrm{ph}}=
\begin{bmatrix}
0.2564{+}\jj 0.7951 & 0.0688{+}\jj 0.3923 & 0.0698{+}\jj 0.3443\\
0.0688{+}\jj 0.3923 & 0.2607{+}\jj 0.7935 & 0.0720{+}\jj 0.3771\\
0.0698{+}\jj 0.3443 & 0.0720{+}\jj 0.3771 & 0.2627{+}\jj 0.7928
\end{bmatrix}\;[\si{\ohm/\kilo\metre}].
\end{equation*}
Separating real and imaginary parts according to \eqref{eq:RLCG}:
\begin{equation*}
\matR=
\begin{bmatrix}
0.2564 & 0.0688 & 0.0698\\
0.0688 & 0.2607 & 0.0720\\
0.0698 & 0.0720 & 0.2627
\end{bmatrix}\!\si{\ohm/\kilo\metre},\;\;
\matL=
\begin{bmatrix}
2.531 & 1.249 & 1.096\\
1.249 & 2.526 & 1.200\\
1.096 & 1.200 & 2.524
\end{bmatrix}\!\si{\milli\henry/\kilo\metre}.
\end{equation*}
Note that the three diagonal inductances differ by less than \SI{0.5}{\percent}
but the three off-diagonals span \(1.10\) to \(1.25\;\si{\milli\henry/\kilo\metre}\):
this asymmetry, induced by the non-symmetric arrangement of phase \(C\)
relative to the \(A\)/\(B\) row, is what drives the eventual current
unbalance.

\subsection{Potential-Coefficient Matrix and Phase Capacitance}

With \(1/(2\pi\varepsilon_{0})\approx 1.7975\times 10^{10}\;\si{\metre/\farad}\), \eqref{eq:Pii}--\eqref{eq:Pij} give, in units of
\(10^{10}\;\si{\metre/\farad}\),
\begin{equation*}
\matP_{\mathrm{full}}=
\begin{bmatrix}
15.191 & 4.148 & 2.935 & 3.513\\
4.148 & 15.191 & 3.875 & 3.950\\
2.935 & 3.875 & 15.443 & 4.272\\
3.513 & 3.950 & 4.272 & 16.366
\end{bmatrix}.
\end{equation*}
Applying \eqref{eq:Kron-P} to eliminate the ground row/column,
\begin{equation*}
\matP_{\mathrm{ph}}=
\begin{bmatrix}
14.437 & 3.300 & 2.018\\
3.300 & 14.238 & 2.844\\
2.018 & 2.844 & 14.328
\end{bmatrix}\;[10^{10}\,\si{\metre/\farad}],
\end{equation*}
and the phase capacitance matrix
\(\matC=\matP_{\mathrm{ph}}^{-1}\) becomes, in \si{\nano\farad/\kilo\metre},
\begin{equation*}
\matC=
\begin{bmatrix}
7.387 & -1.567 & -0.730\\
-1.567 & 7.646 & -1.297\\
-0.730 & -1.297 & 7.341
\end{bmatrix}.
\end{equation*}
Again the entries are visibly non-circulant.

\subsection{\texorpdfstring{Per-Unit-Length \(\matZ\) and \(\matY\)}{Per-Unit-Length Z and Y}}

At \(\omega=100\pi\):
\begin{equation*}
\matZ\;=\;\matR+\jj\omega\matL\cdot 10^{-3},\qquad
\matY\;=\;g_{0}\matI_{3}\cdot 10^{-6}+\jj\omega\matC\cdot 10^{-9},
\end{equation*}
both expressed in \(\si{\ohm/\kilo\metre}\) and \(\si{\siemens/\kilo\metre}\)
respectively. \(\matZ\) coincides with \(\matZ_{\mathrm{ph}}\) above, and
\(\matY\) has diagonal entries \(\approx\jj\,2.32\times 10^{-6}\;\si{\siemens/\kilo\metre}\) with off-diagonals roughly half a microsiemens per kilometre.

\subsection{Sending-End Solution at Selected Line Lengths}

Compute \(\matM=\ee^{\matF\ell}\) numerically, take \(\matA\), \(\matB\) from
its upper blocks, evaluate \(\alpha=\vu\cdot\overline{\matB^{-1}\vu}\) and
\(\beta=\vu\cdot\overline{\matB^{-1}\matA\,\vu}\), and obtain \(W\) from
\eqref{eq:W-closed-form}. The sending current \(\vI_{s}\) follows from
\eqref{eq:Is-from-bc}. The table below collects the result for
\(S=80+\jj 10\;\si{\mega\volt\ampere}\) at four representative lengths:

\begin{center}
\renewcommand{\arraystretch}{1.15}
\begin{tabular}{rrrrrrrr}
\toprule
\(\ell\) [km] & \(|I_{A}|\) [A] & \(|I_{B}|\) [A] & \(|I_{C}|\) [A] &
\(\Delta_{A}\) [\%] & \(\Delta_{B}\) [\%] & \(\Delta_{C}\) [\%] &
\(V_{r}\) [kV LN]\\
\midrule
40   & 407.67 & 446.21 & 417.07 & \(-3.77\) & \(+5.33\) & \(-1.55\) & 59.74 \\
80   & 407.59 & 446.31 & 417.06 & \(-3.79\) & \(+5.35\) & \(-1.56\) & 56.69 \\
120  & 407.51 & 446.41 & 417.05 & \(-3.81\) & \(+5.37\) & \(-1.56\) & 54.49 \\
200  & 407.39 & 446.58 & 417.02 & \(-3.84\) & \(+5.41\) & \(-1.57\) & 53.04 \\
\bottomrule
\end{tabular}
\end{center}
The mean sending current magnitude is
\(\overline{|I|}=\bigl(P^{2}+Q^{2}\bigr)^{1/2}\!/(\sqrt{3}\,V_{LL})\approx
\SI{423.3}{\ampere}\), in good agreement with the row averages above.

\subsection{Interpretation}

\begin{itemize}
\item Phase \(B\) carries roughly \SI{5}{\percent} more current than the
three-phase mean, while phase \(A\) carries \SI{3.8}{\percent} less and
phase \(C\) about \SI{1.6}{\percent} less. The middle phase being the most
loaded is a direct consequence of the geometry: phase \(B\) sits at the
lowest height, equidistant from \(A\) and \(C\), and therefore carries the
largest share of the imbalanced charging current dictated by the non-circulant
\(\matC\) and \(\matL\) matrices computed above.
\item The percentage deviations are remarkably stable over the range
\(40\text{--}\SI{200}{\kilo\metre}\): the structural asymmetry of \(\matZ\)
and \(\matY\) sets the unbalance pattern, not the line length. Length affects
mainly the absolute magnitudes (slight bias because the capacitive charging
current adds to the loading current) and the receiving-end voltage and angle.
\item The receiving-end voltage drops monotonically with \(\ell\), from
\SI{59.7}{\kilo\volt} at \SI{40}{\kilo\metre} (a \SI{5.9}{\percent} drop
relative to \(V_{s}=\SI{63.5}{\kilo\volt}\)) down to \SI{53.0}{\kilo\volt} at
\SI{200}{\kilo\metre} (\SI{16.5}{\percent} drop), reflecting both active
losses across \(\matR\) and reactive drop across \(\matL\).
\item Sweeping \(S\) at fixed \(\ell\) reveals the second axis of the
\(\Delta_{\mathrm{mag},k}(\ell,S)\) surface: as \(|S|\) decreases the
capacitive term \(V_{s}^{2}\beta\) in \eqref{eq:W-closed-form} dominates over
the loading term \(S\), and the percentage unbalance grows in magnitude
because the fixed shunt-induced charging current becomes a larger fraction
of \(|\vI_{s}|\).
\end{itemize}

The full surface
\(\Delta_{\mathrm{mag},k}(\ell,S)\) and the spatial-variation indicator
\(\eta_{k}(\ell,S)\) of \eqref{eq:eta} are produced by the Litgrid
application by repeating the procedure above on a parameter grid.

\end{document}
